\title{: Advancing the Boundaries of Mathematical Reasoning in Open Language Models}

\begin{abstract}
Mathematical reasoning remains a formidable challenge for language models because of its intricate and structured characteristics. In this work, we present DeepSeekMath~7B, which further pre-trains DeepSeek-Coder-Base-v1.5 7B using 120B math-related tokens drawn from Common Crawl, alongside natural language and code datasets. DeepSeekMath~7B attains a notable score of 51.7\% on the competitive MATH benchmark without the use of external toolkits or ensemble voting methods, nearing the performance of Gemini-Ultra and GPT-4. Employing self-consistency across 64 samples from DeepSeekMath~7B yields 60.9\% on MATH. The enhanced mathematical reasoning prowess of DeepSeekMath~can be attributed to two main factors: firstly, leveraging the vast potential of publicly available web data through a carefully designed data selection process; secondly, introducing Group Relative Policy Optimization (GRPO), a PPO variant that boosts mathematical reasoning capabilities while simultaneously optimizing PPO’s memory efficiency.
\end{abstract}

\section{Introduction}

Large language models (LLMs) have transformed approaches to mathematical reasoning within artificial intelligence, driving major progress on both quantitative reasoning benchmarks \citep{MATH} and geometric reasoning benchmarks \citep{trinh2024solving}. Furthermore, these models have demonstrated utility in aiding humans to tackle complex mathematical problems \citep{tao}. Nonetheless, state-of-the-art models such as GPT-4 \citep{gpt4} and Gemini-Ultra \citep{gemini} are not openly accessible, and currently available open-source alternatives lag significantly behind in terms of performance.

In this work, we present DeepSeekMath, a language model tailored for the mathematical domain that considerably surpasses the mathematical performance of existing open-source models and nears the proficiency of GPT-4 on academic evaluation tasks. To accomplish this, we develop the DeepSeekMath Corpus, a large-scale, high-quality pre-training dataset consisting of 120B math tokens. This corpus is derived from the Common Crawl (CC) through the use of a fastText-based classifier \citep{joulin2016fasttext}. Initially, the classifier is trained with positive samples drawn from OpenWebMath \citep{openwebmath}, alongside a varied collection of other web pages as negative samples. We then utilize this classifier to extract additional positive samples from the CC, which are subsequently refined via human annotation. The classifier is updated with these enhanced data to boost its accuracy. Evaluation results demonstrate that the large-scale corpus is of superior quality, as evidenced by our base model DeepSeekMath-Base 7B achieving 64.2\% on GSM8K \citep{gsm8k} and 36.2\% on the challenging MATH dataset \citep{MATH}, outperforming Minerva 540B \citep{minerva}. Moreover, since the DeepSeekMath Corpus is multilingual, we observe improvements on Chinese mathematical benchmarks \citep{wei2023cmath,agieval}. We consider our experience in processing mathematical data as an initial contribution to the research community, with ample opportunities for future enhancements.

DeepSeekMath-Base is initialized from DeepSeek-Coder-Base-v1.5 7B \citep{deepseek-coder}, as we find that beginning with a model trained on code is a superior strategy compared to using a general LLM. Additionally, we observe that the mathematical training also enhances the model's performance on MMLU \citep{mmlu} and BBH benchmarks \citep{bbh}, suggesting that it not only improves mathematical proficiency but also strengthens general reasoning skills.

Following pre-training, we perform mathematical instruction tuning on DeepSeekMath-Base utilizing data involving chain-of-thought \citep{cot}, program-of-thought \citep{pot,pal}, and tool-integrated reasoning \citep{tora}. The resulting model, DeepSeekMath-Instruct 7B, outperforms all other 7B-scale models and achieves performance comparable to open-source instruction-tuned models at the 70B scale.

Moreover, we propose the Group Relative Policy Optimization (GRPO), a variant of the Proximal Policy Optimization (PPO) reinforcement learning algorithm \citep{schulman2017proximal}. GRPO eliminates the use of a critic model and instead derives the baseline from group-level scores, substantially reducing the computational resources needed for training. Utilizing only a subset of English instruction tuning data, GRPO achieves significant gains over the robust DeepSeekMath-Instruct model across both in-domain tasks (GSM8K: 82.9\% $\rightarrow$ 88.2\%, MATH: 46.8\% $\rightarrow$ 51.7\%) and out-of-domain mathematical tasks (e.g., CMATH: 84.6\% $\rightarrow$ 88.8\%) during the reinforcement learning stage. We also offer a unified framework to interpret various approaches such as Rejection Sampling Fine-Tuning (RFT) \citep{yuan2023scaling}, Direct Preference Optimization (DPO) \citep{dpo}, PPO, and GRPO. Within this framework, we recognize that these methods can be viewed as either direct or simplified reinforcement learning techniques. We conduct comprehensive experiments—such as comparing online versus offline training, outcome versus process supervision, and single-turn versus iterative RL—to thoroughly explore the fundamental aspects of this framework. Finally, we clarify why our RL approach enhances the performance of instruction-tuned models and outline prospective directions for developing more effective RL methods based on this unified framework.

\subsection{Contributions}  
Our work presents scalable mathematical pre-training, as well as an investigation and evaluation of reinforcement learning techniques.

\textbf{Math Pre-Training at Scale}

\begin{itemize}[topsep=0pt]
    \item We demonstrate strong evidence that publicly available Common Crawl data holds valuable mathematical information. Through a carefully crafted data selection process, we create the DeepSeekMath Corpus, a high-quality dataset comprising 120B tokens derived from math-focused web pages. This dataset is nearly 7 times larger than the math web pages utilized by Minerva \citep{minerva} and 9 times bigger than the newly released OpenWebMath \citep{openwebmath}.

    \item Our pre-trained base model, DeepSeekMath-Base 7B, attains performance on par with Minerva 540B \citep{minerva}, suggesting that the number of parameters alone does not determine mathematical reasoning ability. Smaller models trained on high-quality datasets can also achieve strong results.

    \item We report insights from our math training experiments. Pre-training on code before math training enhances models' proficiency in solving mathematical problems, both with and without the use of external tools. This provides a partial response to the longstanding question: \textit{does code training enhance reasoning skills?} Our results indicate that it does, at least regarding mathematical reasoning.

    \item Although training on arXiv papers is widespread, particularly in many math-related studies, it does not yield significant improvements across all mathematical benchmarks evaluated in this work.
\end{itemize}

\textbf{Exploration and Analysis of Reinforcement Learning}

\begin{itemize}[topsep=0pt]
    \item We propose Group Relative Policy Optimization (GRPO), a reinforcement learning algorithm that is both efficient and effective. Unlike Proximal Policy Optimization (PPO), GRPO eliminates the need for a critic model by estimating the baseline using group scores, thereby significantly lowering the training resource requirements.
    \item We show that GRPO markedly improves the performance of our instruction-tuned model DeepSeekMath-Instruct, relying solely on instruction-tuning data. Additionally, we observe gains in out-of-domain performance throughout the reinforcement learning process.
    \item We offer a unified framework to interpret various methods such as RFT, DPO, PPO, and GRPO. Furthermore, we perform extensive experiments comparing online versus offline training, outcome versus process supervision, single-turn versus iterative reinforcement learning, and more, to thoroughly explore the key components of this framework.
    \item Drawing from our unified paradigm, we investigate the underlying reasons for the success of reinforcement learning and outline several promising avenues for achieving more effective reinforcement learning in large language models (LLMs).
\end{itemize}

\subsection{Summary of Evaluations and Metrics}

\begin{itemize}[topsep=0pt]
    \item \textbf{English and Chinese Mathematical Reasoning}: We carry out thorough evaluations of our models on both English and Chinese benchmarks, which include math problems ranging from elementary to university level. The English datasets comprise GSM8K \citep{gsm8k}, MATH \citep{MATH}, SAT \citep{llemma}, OCW Courses \citep{minerva}, and MMLU-STEM \citep{mmlu}. The Chinese datasets include MGSM-zh \citep{mgsm}, CMATH \citep{wei2023cmath}, Gaokao-MathCloze \citep{agieval}, and Gaokao-MathQA \citep{agieval}. We assess the models’ ability to produce self-contained textual solutions without utilizing tools, as well as their capability to solve problems using Python.

    On English benchmarks, DeepSeekMath-Base competes well against the closed-source Minerva 540B \citep{minerva} and outperforms all open-source base models (e.g., Mistral 7B \citep{mistral} and Llemma-34B \citep{llemma}), whether or not these models have undergone math pre-training, often by a wide margin. Importantly, DeepSeekMath-Base achieves superior results on Chinese benchmarks, likely because we do not restrict our math pre-training data to English only as previous works \citep{minerva,llemma} have done, and also incorporate high-quality non-English datasets. Following mathematical instruction tuning and reinforcement learning, our DeepSeekMath-Instruct and DeepSeekMath-RL models exhibit robust performance, achieving over 50\% accuracy on the competitive MATH dataset—the first such accomplishment in the open-source community.

\begin{itemize}
    \item \textbf{Formal Mathematics}: We assess DeepSeekMath-Base on the informal-to-formal theorem proving task introduced by \citep{dsp_proof} using the miniF2F dataset \citep{minif2f}, with Isabelle \citep{isabelle} selected as the proof assistant. DeepSeekMath-Base exhibits strong few-shot autoformalization capabilities.
    \item \textbf{Natural Language Understanding, Reasoning, and Code}: To obtain a thorough evaluation of the models' overall comprehension, reasoning, and coding skills, we test DeepSeekMath-Base on the Massive Multitask Language Understanding (MMLU) benchmark \citep{mmlu}, which includes 57 multiple-choice tasks spanning various disciplines; BIG-Bench Hard (BBH) \citep{bbh}, composed of 23 demanding tasks that primarily require multi-step reasoning; as well as HumanEval \citep{codex} and MBPP \citep{mbpp}, both commonly used for assessing code language models. Math pre-training positively influences performance in language understanding and reasoning.
\end{itemize}

\section{Math Pre-Training}

\subsection{Data Collection and Decontamination} In this section, we describe the methodology for creating the DeepSeekMath Corpus from Common Crawl. As illustrated in Figure \ref{fig:data_collect}, we propose an iterative pipeline that methodically collects a large-scale mathematical corpus from Common Crawl, starting with a seed dataset (for instance, a relatively small but high-quality math-related dataset). This approach is also adaptable to other fields, such as programming.

Initially, we select OpenWebMath \citep{openwebmath}, a curated collection of high-quality mathematical web content, as our seed corpus. Using this dataset, we train a fastText model \citep{joulin2016fasttext} to retrieve additional web pages resembling those in OpenWebMath. Concretely, we randomly pick 500,000 samples from the seed corpus as positive examples and another 500,000 web pages from Common Crawl as negative examples. The model is trained using an open-source library\footnote{\url{https://fasttext.cc}}, configured with a vector size of 256, a learning rate of 0.1, a maximum word n-gram length of 3, a minimum word count of 3, and 3 training epochs. To reduce the scale of the original Common Crawl, we apply URL-based deduplication and near-duplicate removal, yielding 40 billion HTML web pages. We then use the fastText model to retrieve mathematical pages from this deduplicated subset. To eliminate low-quality math content, we rank the retrieved pages based on their fastText scores and retain only the highest-scoring ones. The amount of data retained is evaluated through pre-training experiments on top 40B, 80B, 120B, and 160B tokens. In the first iteration, we opt to keep the top 40 billion tokens.

Following the initial round of data collection, many mathematical web pages remain uncollected, primarily because the fastText model is trained on a set of positive samples that lacks adequate diversity. To address this, we identify additional mathematical web sources to expand the seed corpus, allowing us to further optimize the fastText model. Specifically, we first partition the entire Common Crawl into separate domains, where a domain is defined as web pages sharing the same base URL. For each domain, we compute the proportion of web pages collected during the first iteration. Domains for which more than 10\% of the web pages have been gathered are labeled as math-related (e.g., \url{mathoverflow.net}).

Next, we manually annotate URLs containing mathematical content within these identified domains (e.g., \url{mathoverflow.net/questions}). Any web pages linked to these URLs that were not collected previously are then incorporated into the seed corpus. This strategy allows us to accumulate additional positive examples, thus training an enhanced fastText model capable of retrieving more mathematical data in the next iteration. After performing four iterations of data collection, we accumulate 35.5M mathematical web pages, amounting to 120B tokens. By the fourth iteration, we observe that almost 98\% of the data had already been collected by the third iteration, leading us to halt further data collection.

To prevent contamination of benchmarks, we adopt the approach of \cite{deepseek-coder} by filtering out web pages containing questions or answers from English mathematical benchmarks such as GSM8K \citep{gsm8k} and MATH \citep{MATH}, as well as Chinese benchmarks like CMATH \citep{wei2023cmath} and AGIEval \citep{agieval}. The filtering process removes any text segment that contains a 10-gram sequence matching exactly any substring from these evaluation benchmarks. For benchmark texts shorter than 10 grams but at least 3 grams in length, we apply exact matching to exclude contaminated web pages from our math training corpus.

\subsection{Evaluating the Quality of the DeepSeekMath~Corpus}

We conduct pre-training experiments to compare the DeepSeekMath~Corpus with other recently published math-focused training corpora:

\begin{itemize}[topsep=0pt]
    \item \textbf{MathPile} \citep{mathpile}: a multi-source dataset (8.9B tokens) compiled from textbooks, Wikipedia, ProofWiki, CommonCrawl, StackExchange, and arXiv, with over 85\% of the content originating from arXiv;
    \item \textbf{OpenWebMath} \citep{openwebmath}: CommonCrawl data filtered specifically for mathematical content, totaling 13.6B tokens;
    \item \textbf{Proof-Pile-2} \citep{llemma}: a mathematical dataset comprising OpenWebMath, AlgebraicStack (10.3B tokens of mathematical code), and arXiv papers (28.0B tokens). For experiments using Proof-Pile-2, we follow \cite{llemma} in employing an arXiv:Web:Code token ratio of 2:4:1.
\end{itemize}

\subsubsection{Training Setup}
\label{sec:quality-policy}
We fine-tune a general pre-trained language model containing 1.3 billion parameters, utilizing the same architecture as the DeepSeek LLMs \citep{deepseek-llm}, referred to as DeepSeek-LLM 1.3B. Each mathematical dataset is used independently to train the model for 150 billion tokens. All training experiments employ the efficient and lightweight HAI-LLM \citep{haillm} framework. Adhering to the DeepSeek LLMs' training protocol, we adopt the AdamW optimizer \citep{adamW} with parameters $\beta_1=0.9$, $\beta_2=0.95$, and a weight decay of 0.1. The learning rate schedule involves multiple phases: it gradually increases to its maximum over 2,000 warm-up steps, then decays to 31.6\% of the peak after 80\% of the training duration, and further drops to 10.0\% of the peak following 90\% of the training time. The peak learning rate is set to 5.3e-4, with a batch size of 4 million tokens and a context length of 4,000 tokens.

\subsubsection{Evaluation Outcomes}

\textbf{The DeepSeekMath~Corpus demonstrates superior quality, supports multiple languages for mathematical content, and is the largest dataset available.}

\begin{itemize}[topsep=0pt]
    \item \textbf{High quality}:
    We assess downstream task performance on eight mathematical benchmarks using few-shot chain-of-thought prompting \cite{cot}. As presented in Table \ref{tab:corpora_comparison}, models trained on the DeepSeekMath~Corpus consistently outperform others. Figure \ref{fig:corpora_comparisons} further illustrates that at 50 billion tokens (equivalent to one full epoch over Proof-Pile-2), the model trained on DeepSeekMath surpasses Proof-Pile-2, highlighting the superior average quality of the DeepSeekMath dataset.
    \item \textbf{Multilingual support}:
    The DeepSeekMath~Corpus contains data spanning several languages, with English and Chinese being the predominant ones. According to Table \ref{tab:corpora_comparison}, training on DeepSeekMath improves mathematical reasoning capabilities in both English and Chinese. Conversely, current mathematical datasets, which are mainly English-focused, provide limited enhancement and may even degrade performance for Chinese mathematical reasoning.
    \item \textbf{Large scale}:
    The DeepSeekMath~Corpus exceeds existing mathematical datasets in size by several folds. As shown in Figure \ref{fig:corpora_comparisons}, the DeepSeek-LLM 1.3B trained on DeepSeekMath exhibits a steeper and more sustained learning curve. In contrast, smaller baseline corpora have undergone multiple reuses during training, leading the model's performance to quickly plateau.
\end{itemize}

\subsection{Training and Evaluating DeepSeekMath-Base 7B}

In this section, we present DeepSeekMath-Base 7B, a foundational model exhibiting strong reasoning capabilities, particularly in mathematics. The model is initialized from DeepSeek-Coder-Base-v1.5 7B \citep{deepseek-coder} and trained on 500 billion tokens. The training data composition is as follows: 56\% originates from the DeepSeekMath Corpus, 4\% from AlgebraicStack, 10\% from arXiv, 20\% consists of Github code, and the remaining 10\% comprises natural language data drawn from Common Crawl in both English and Chinese. We largely follow the training protocol described in Section \ref{sec:quality-policy}, except that we set the peak learning rate to 4.2e-4 and employ a batch size of 10 million tokens.

We perform an extensive evaluation of DeepSeekMath-Base 7B’s mathematical abilities, emphasizing its capacity to generate self-contained mathematical solutions without external tool use, solve math problems via tools, and execute formal theorem proving. In addition to its mathematical prowess, we provide a broader assessment of the base model’s overall performance in natural language understanding, reasoning, and programming skills.

\paragraph{Mathematical Problem Solving with Step-by-Step Reasoning}

We test DeepSeekMath-Base’s ability to solve mathematical problems using few-shot chain-of-thought prompting \citep{cot} across eight benchmarks in both English and Chinese. These benchmarks cover quantitative reasoning tasks (such as GSM8K \citep{gsm8k}, MATH \citep{MATH}, and CMATH \citep{wei2023cmath}) and multiple-choice problems (including MMLU-STEM \citep{mmlu} and Gaokao-MathQA \citep{agieval}), representing a wide range of mathematical domains from elementary to university-level difficulty.

As indicated in Table \ref{tab:base_math_cot}, DeepSeekMath-Base 7B achieves the highest performance among open-source base models across all eight benchmarks. This group includes the popular general-purpose model Mistral 7B \citep{mistral} and the recently introduced Llemma 34B \citep{llemma}, which was trained on mathematical data from Proof-Pile-2 \citep{llemma}. Significantly, on the competition-level MATH dataset, DeepSeekMath-Base outperforms all existing open-source base models by more than 10\% in absolute terms and even surpasses Minerva 540B \citep{minerva}, a closed-source base model that is 77 times larger, based on PaLM \citep{palm} and further trained on mathematical texts.

\paragraph{Mathematical Problem Solving with Tool Use} We assess program-assisted mathematical reasoning on GSM8K and MATH using few-shot program-of-thought prompting \citep{pot,pal}. Models are tasked with solving each problem by generating a Python program, leveraging libraries such as \textit{math} and \textit{sympy} for complex calculations. The output from running the program serves as the answer. As presented in Table \ref{tab:base_math_pot_proof}, DeepSeekMath-Base 7B surpasses the previous state-of-the-art Llemma 34B.

\paragraph{Formal Mathematics}

Automating formal proof generation is valuable for verifying the correctness and reliability of mathematical proofs and improving efficiency, a focus of growing interest recently. We test DeepSeekMath-Base 7B on the informal-to-formal proof task from \citep{dsp_proof}, which involves creating a formal proof given an informal statement, its formal equivalent, and an informal proof. Using miniF2F \citep{minif2f}, a benchmark for formal Olympiad-level mathematics, we generate formal proofs in Isabelle for each problem via few-shot prompting. Following \cite{dsp_proof}, models first produce proof sketches, after which the off-the-shelf automated prover Sledgehammer \citep{sledgehammer} completes the missing details. As shown in Table \ref{tab:base_math_pot_proof}, DeepSeekMath-Base 7B demonstrates robust results in automatic proof formalization.

\paragraph{Natural Language Understanding, Reasoning, and Code}

We evaluate model capabilities in natural language understanding using MMLU \citep{mmlu}, reasoning with BBH \citep{bbh}, and coding through HumanEval \citep{codex} and MBPP \citep{mbpp}. As illustrated in Table \ref{tab:understanding_reasoning_code}, DeepSeekMath-Base 7B shows marked improvements over its predecessor, DeepSeek-Coder-Base-v1.5 \citep{deepseek-coder}, on MMLU and BBH, demonstrating the beneficial effect of math training on language comprehension and reasoning skills. Furthermore, by incorporating code tokens during continued training, DeepSeekMath-Base 7B maintains the coding performance level of DeepSeek-Coder-Base-v1.5 on the two code benchmarks. Overall, DeepSeekMath-Base 7B considerably outperforms the general-purpose Mistral 7B \citep{mistral} on all three reasoning and coding benchmarks.

\section{Supervised Fine-Tuning}

\subsection{SFT Data Preparation}

We compile a mathematical instruction-tuning dataset encompassing English and Chinese problems from various mathematical domains and with different difficulty levels: each problem is paired with solutions presented in chain-of-thought (CoT) \citep{cot}, program-of-thought (PoT) \citep{pot,pal}, and tool-assisted reasoning formats \citep{tora}. The dataset consists of a total of 776K training instances.

\begin{itemize}[topsep=0pt]
    \item \textbf{English mathematical datasets}: We annotate GSM8K and MATH problems with tool-assisted solutions, and utilize a subset of MathInstruct \citep{MathInstruct} along with the Lila-OOD training set \citep{lila}, where problems are solved using CoT or PoT. Our English dataset covers a wide range of mathematical fields, including algebra, probability, number theory, calculus, and geometry.
    \item \textbf{Chinese mathematical datasets}: We gather Chinese K-12 math problems covering 76 sub-topics such as linear equations, with solutions annotated in both CoT and tool-assisted reasoning formats.
\end{itemize}

\subsection{Training and Evaluation of DeepSeekMath-Instruct 7B}

Here, we present DeepSeekMath-Instruct 7B, which undergoes mathematical instruction tuning built upon DeepSeekMath-Base. Training samples are concatenated randomly until the maximum context length of 4K tokens is reached. The model is trained for 500 steps with a batch size of 256 and a fixed learning rate of 5e-5.

We assess the mathematical capabilities of the models both without and with tool usage across four quantitative reasoning benchmarks in English and Chinese. We compare our model against state-of-the-art models at the time:

\begin{itemize}[topsep=0pt]
    \item \textbf{Closed-source models} include: (1) the GPT series, where GPT-4 \citep{gpt4} and GPT-4 Code Interpreter~\footnote{\url{https://openai.com/blog/chatgpt-plugins\#\#code-interpreter}} are the most advanced, (2) Gemini Ultra and Pro \citep{gemini}, (3) Inflection-2 \citep{inflection-2}, (4) Grok-1~\footnote{\url{https://x.ai/model-card}}, as well as recently released models by Chinese companies, including (5) Baichuan-3~\footnote{\url{https://www.baichuan-ai.com}}, and (6) the newest GLM-4~\footnote{\url{https://open.bigmodel.cn/dev/api\#glm-4}} from the GLM family \citep{glm}.

These models serve general purposes, most of which have undergone a series of alignment procedures.
\item \textbf{Open-source models} consist of: general-purpose models such as (1) DeepSeek-LLM-Chat 67B \citep{deepseek-llm}, (2) Qwen 72B \citep{qwen}, (3) SeaLLM-v2 7B \citep{seallm}, and (4) ChatGLM3 6B \citep{chatglm3}, along with models enhanced in mathematical capabilities including (5) InternLM2-Math 20B~\footnote{\url{https://github.com/InternLM/InternLM-Math}}, which is based on InternLM2 and was trained on math tasks followed by instruction tuning, (6) Math-Shepherd-Mistral 7B that employs PPO training \citep{schulman2017proximal} applied to Mistral 7B \citep{mistral} using a process-supervised reward model, (7) the WizardMath series \citep{wizardmath}, which enhances mathematical reasoning in Mistral 7B and Llama-2 70B \citep{llama2} through evolve-instruct (a variant of instruction tuning using AI-generated instructions) and PPO training, with training problems mainly derived from GSM8K and MATH datasets, (8) MetaMath 70B \citep{metamath}, a fine-tuned Llama-2 70B on an augmented version of GSM8K and MATH, (9) ToRA 34B \cite{tora}, which fine-tunes CodeLlama 34B for tool-integrated mathematical reasoning, and (10) MAmmoTH 70B \citep{MathInstruct}, an instruction-tuned Llama-2 70B based on MathInstruct.
\end{itemize}

As indicated in Table \ref{tab:sft_rl_math}, in an evaluation scenario where tool usage is prohibited, DeepSeekMath-Instruct 7B exhibits strong step-by-step reasoning capabilities. Remarkably, on the competition-level MATH dataset, our model outperforms all open-source models and most proprietary models (e.g., Inflection-2 and Gemini Pro) by a margin of at least 9\% absolute. This advantage holds even against considerably larger models (e.g., Qwen 72B) or those specifically improved through math-centered reinforcement learning (e.g., WizardMath-v1.1 7B). Although DeepSeekMath-Instruct matches Chinese proprietary models GLM-4 and Baichuan-3 on MATH, it still falls short compared to GPT-4 and Gemini Ultra.

In the evaluation setting allowing the combination of natural language reasoning and program-based tool usage for solving problems, DeepSeekMath-Instruct 7B approaches 60\% accuracy on MATH, exceeding all current open-source models. Across other benchmarks, our model performs competitively with DeepSeek-LLM-Chat 67B, the previous state-of-the-art model that is ten times larger.

\section{Reinforcement Learning}

\subsection{Group Relative Policy Optimization} Reinforcement learning (RL) has been demonstrated to be effective in further enhancing the mathematical reasoning capabilities of LLMs after the Supervised Fine-Tuning (SFT) phase \citep{wang2023math,wizardmath}. Here, we present our efficient and effective RL algorithm, Group Relative Policy Optimization (GRPO).

\subsubsection{From PPO to GRPO} Proximal Policy Optimization (PPO) \citep{schulman2017proximal} is an actor-critic RL method widely employed during RL fine-tuning of LLMs \citep{ouyang2022training}. Specifically, it updates LLMs by maximizing the following surrogate objective:

\begin{equation}
\footnotesize
    \mathcal{J}_{PPO}(\theta) = \mathbb{E}_{q \sim P(Q), o \sim \pi_{\theta_{old}}(O|q)} \frac{1}{|o|} \sum_{t=1}^{|o|} \min \left[ \frac{\pi_\theta(o_{t} | q, o_{<t})}{\pi_{\theta_{old}}(o_{t} | q, o_{<t})} A_{t}, \text{clip} \left( \frac{\pi_\theta(o_{t} | q, o_{<t})}{\pi_{\theta_{old}}(o_{t} | q, o_{<t})}, 1 - \varepsilon, 1 + \varepsilon \right)  A_{t} \right]

Here, $\pi_{\theta}$ and $\pi_{\theta_{old}}$ denote the current and previous policy models, while $q$ and $o$ represent questions and outputs sampled from the question dataset and the former policy $\pi_{\theta_{old}}$, respectively. The parameter $\varepsilon$ is a clipping-related hyper-parameter introduced in PPO to ensure training stability. The term $A_t$ refers to the advantage, which is computed using Generalized Advantage Estimation (GAE) \citep{gae}, relying on the sequence of rewards $\{r_{\ge t}\}$ and a learned value function $V_{\psi}$. Consequently, PPO requires simultaneous training of a value function alongside the policy model. To prevent excessive optimization of the reward model, the standard technique is to incorporate a per-token KL penalty from a reference model into the reward at each token \citep{ouyang2022training}, formally expressed as

\begin{equation}
    r_{t} = r_\varphi(q, o_{\le t}) - \beta \log\frac{\pi_{\theta}(o_{t}|q, o_{<t})}{\pi_{ref}(o_{t}|q, o_{<t})},
\label{eq:PPO-reward}
\end{equation}

where $r_\varphi$ is the reward model, $\pi_{ref}$ is the reference model—commonly the initial SFT model—and $\beta$ is the coefficient controlling the KL penalty.

Since the value function in PPO is typically another model of size comparable to the policy model, it introduces significant memory and computational costs. Moreover, during RL training, the value function acts as a baseline in computing the advantage to reduce variance. In the context of large language models (LLMs), usually only the final token receives a reward from the reward model, which complicates training an accurate value function at every token. To overcome this challenge, as illustrated in Figure \ref{fig:grpo}, we introduce Group Relative Policy Optimization (GRPO), which eliminates the need for an additional value function approximation as required by PPO. Instead, GRPO uses the average reward of multiple sampled outputs generated in response to the same question as the baseline. Specifically, for each question $q$, GRPO draws a group of outputs $\{o_1, o_2, \cdots, o_G\}$ from the old policy $\pi_{\theta_{old}}$, then optimizes the policy by maximizing the following objective:

\begin{equation}
\footnotesize
\begin{split}
    \mathcal{J}_{GRPO}(\theta) &= \mathbb{E}{[q \sim P(Q), \{o_i\}_{i=1}^G \sim \pi_{\theta_{old}}(O|q)]}  \\
    & \frac{1}{G}\sum_{i=1}^G\frac{1}{|o_i|} \sum_{t=1}^{|o_i|} \left\{ \min \left[ \frac{\pi_\theta(o_{i,t} | q, o_{i,<t})}{\pi_{\theta_{old}}(o_{i,t} | q, o_{i,<t})} \hat{A}_{i,t}, \text{clip} \left( \frac{\pi_\theta(o_{i,t} | q, o_{i,<t})}{\pi_{\theta_{old}}(o_{i,t} | q, o_{i,<t})}, 1 - \varepsilon, 1 + \varepsilon \right)  \hat{A}_{i,t} \right] - \beta \mathbb{D}_{KL}\left[\pi_{\theta} || \pi_{ref}\right]\right\} ,
\end{split}
\label{eq:GRPO-obj}
\end{equation}

where $\varepsilon$ and $\beta$ are hyper-parameters, and $\hat{A}_{i,t}$ denotes the advantage calculated solely based on the relative rewards of outputs within each group, as detailed in subsequent sections. The group-relative strategy employed by GRPO to compute advantages aligns naturally with the comparative characteristic of reward models, since these models are generally trained on datasets containing comparisons between outputs for the same question. It is also noteworthy that, instead of incorporating the KL penalty into the reward as in (\ref{eq:PPO-reward}), GRPO applies regularization by adding the KL divergence between the trained policy and the reference policy directly to the loss, thus simplifying the calculation of $\hat{A}_{i,t}$. Distinct from the KL penalty term in (\ref{eq:PPO-reward}), the KL divergence here is estimated using the following unbiased estimator \citep{kl_approx}:

\begin{algorithm}[t]
  \small
  \caption{Iterative Group Relative Policy Optimization}
  \textbf{Input} initial policy model $\pi_{\theta_{\text{init}}}$; reward models $r_\varphi$; task prompts $\mathcal{D}$; 
  hyperparameters $\varepsilon$, $\beta$, $\mu$
  \begin{algorithmic}[1]
    \State Initialize policy model $\pi_\theta \leftarrow \pi_{\theta_{\text{init}}}$
    \For{iteration = 1, \dots, I}
       \State Set reference model $\pi_{ref} \leftarrow \pi_{\theta}$
      \For{step = 1, \dots, M}
      \State Draw a batch $\mathcal{D}_b$ from $\mathcal{D}$
      \State Store the current policy as $\pi_{\theta_{old}} \leftarrow \pi_{\theta}$ 
      \State For each question $q \in \mathcal{D}_b$, sample $G$ outputs $\{o_i\}_{i=1}^G \sim \pi_{\theta_{old}} (\cdot \mid q) $
      \State Evaluate each sampled output $o_i$ using $r_{\varphi}$ to obtain rewards $\{r_i\}_{i=1}^{G}$
      \State Calculate $\hat{A}_{i,t}$ for the $t$-th token of $o_i$ via group relative advantage estimation.
      \For{GRPO iteration = 1, \dots, $\mu$}
        \State Update the policy $\pi_{\theta}$ by maximizing the GRPO objective (Equation \ref{eq:GC-GRPO})
      \EndFor
    \EndFor \State Update reward model $r_\varphi$ through continual training with a replay buffer.
    \EndFor 
  \end{algorithmic}
  \textbf{Output} $\pi_\theta$
  \label{alg:iter-grpo}
\end{algorithm}

\subsubsection{Outcome Supervision RL with GRPO} Formally, for each query $q$, a set of outputs $\{o_1, o_2, \ldots, o_G\}$ is generated by sampling from the previous policy $\pi_{\theta_{old}}$. These outputs are then scored by a reward model, producing a reward vector $\mathbf{r} = \{r_1, r_2, \ldots, r_G\}$. Next, the rewards are normalized by subtracting their group mean and dividing by the group standard deviation. Outcome supervision assigns this normalized reward at the conclusion of each output $o_i$ and sets the advantage values $\hat{A}_{i,t}$ for all tokens in the output equal to this normalized reward, i.e., $\hat{A}_{i,t} = \widetilde{r}_i = \frac{r_i - {\rm mean}(\mathbf{r})}{{\rm std}(\mathbf{r})}$. The policy is then optimized by maximizing the objective described in equation (\ref{eq:GRPO-obj}).

\subsubsection{Process Supervision RL with GRPO} Outcome supervision provides only a terminal reward for each output, which might be insufficient or inefficient for guiding the policy on complex mathematical problems. Inspired by \cite{wang2023math}, we also investigate process supervision, which supplies rewards at every reasoning step. Specifically, given question $q$ and $G$ sampled outputs $\{o_1, o_2, \ldots, o_G\}$, a process reward model evaluates each reasoning step within the outputs, producing rewards: $\mathbf{R} = \{\{r_1^{index(1)}, \ldots, r_1^{index(K_1)}\}, \ldots, \{r_G^{index(1)}, \ldots, r_G^{index(K_G)}\}\}$, where $index(j)$ denotes the token index marking the end of the $j$-th step, and $K_i$ is the total number of steps in output $i$. These step-wise rewards are also normalized using the overall mean and standard deviation as $\widetilde{r}_i^{index(j)} = \frac{r_i^{index(j)} - {\rm mean}(\mathbf{R})}{{\rm std}(\mathbf{R})}$. The advantage for each token is then computed as the sum of normalized rewards from all subsequent steps, i.e., $\hat{A}_{i,t} = \sum_{index(j) \ge t} \widetilde{r}_i^{index(j)}$, and the policy is optimized by maximizing the objective in equation (\ref{

\subsubsection{Iterative RL with GRPO} As the reinforcement learning training advances, the initial reward model may become inadequate for guiding the current policy model. Hence, we investigate iterative RL using GRPO. As illustrated in Algorithm \ref{alg:iter-grpo}, in iterative GRPO, new training datasets for the reward model are created from samples generated by the policy model. The previous reward model is then continuously updated using a replay strategy that incorporates 10\% of past data. After that, the policy model is set as the reference model and is further trained using the updated reward model.

\subsection{Training and Evaluating DeepSeekMath-RL}

We perform RL training based on DeepSeekMath-Instruct 7B. The RL training data consist of chain-of-thought formatted questions related to GSM8K and MATH drawn from the SFT dataset, which contains approximately 144K questions. Other SFT questions are excluded to examine the effects of RL on benchmarks with limited data during the RL phase. The reward model’s training dataset is constructed according to \citep{wang2023math}. Our initial reward model is trained starting from DeepSeekMath-Base 7B with a learning rate of 2e-5. For GRPO, the policy model’s learning rate is set to 1e-6, and the KL coefficient is fixed at 0.04. For each question, we generate $64$ outputs. The maximum sequence length is 1024, and the training batch size is 1024. The policy model undergoes only one update after each exploration phase. DeepSeekMath-RL 7B is evaluated on benchmarks following the same protocol as DeepSeekMath-Instruct 7B. In this context, GSM8K and MATH with chain-of-thought reasoning are considered in-domain tasks, while all other benchmarks are treated as out-of-domain tasks.

Table \ref{tab:sft_rl_math} presents the results of open- and closed-source models employing both chain-of-thought and tool-enhanced reasoning on benchmarks in English and Chinese. Our observations are: 1) DeepSeekMath-RL 7B achieves accuracies of 88.2\% on GSM8K and 51.7\% on MATH when using chain-of-thought reasoning. This performance exceeds that of all open-source models ranging from 7B to 70B parameters, as well as most closed-source models. 2) Importantly, DeepSeekMath-RL 7B is trained solely on chain-of-thought-format instruction tuning data from GSM8K and MATH, building upon DeepSeekMath-Instruct 7B. Despite the limited training dataset, it surpasses DeepSeekMath-Instruct 7B on all evaluation metrics, demonstrating the effectiveness of reinforcement learning.

\section{Discussion}
Here, we discuss our insights derived from pre-training and reinforcement learning experiments.

\subsection{Insights from Pre-Training}

We begin by sharing our experiences related to pre-training. Unless otherwise noted, we follow the training configurations described in Section \ref{sec:quality-policy}. It is important to clarify that when mentioning the DeepSeekMath Corpus in this section, we refer to an 89B-token dataset collected during the second iteration of data gathering.

\subsubsection{Code Training Enhances Mathematical Reasoning}

A widely held but unconfirmed hypothesis posits that training on code improves reasoning skills. We aim to provide partial evidence supporting this claim, specifically in the mathematical context: training on code enhances models’ capacity for mathematical reasoning both with and without tool assistance.

To examine how code training influences mathematical reasoning, we conducted experiments using the following training setups: two-stage training and one-stage training.

\textbf{Two-Stage Training}

\begin{itemize}[topsep=0pt]
    \item \textbf{Code Training for 400B Tokens $\rightarrow$ Math Training for 150B Tokens}: We initially train DeepSeek-LLM 1.3B on 400B code tokens, followed by a subsequent training phase on 150B math tokens;
    \item \textbf{General Training for 400B Tokens $\rightarrow$ Math Training for 150B Tokens}: As a control, we also perform training using general tokens (drawn from a large-scale general corpus compiled by DeepSeek-AI) instead of code tokens during the first phase, aiming to explore the benefits of code tokens compared to general tokens for enhancing mathematical reasoning.
\end{itemize}

\textbf{One-Stage Training}

\begin{itemize}[topsep=0pt]
    \item \textbf{Math Training for 150B Tokens}: We train DeepSeek-LLM 1.3B solely on 150B math tokens;
    \item \textbf{Training on a mixture of 400B Code Tokens and 150B Math Tokens}: Since math training after code training degrades coding ability, we explore whether combining code tokens and math tokens in a single-stage training can still improve mathematical reasoning while also mitigating catastrophic forgetting.
\end{itemize}

\paragraph{Results}
Tables \ref{tab:code-math} and \ref{tab:code-math-general-eval} present downstream task performance across various training configurations.

Training on code tokens enhances program-assisted mathematical reasoning in both two-stage and one-stage training frameworks. As indicated in Table \ref{tab:code-math}, in the two-stage setting, code training alone considerably improves the capacity to solve GSM8K and MATH challenges using Python, with further gains achieved after the subsequent math training stage. Notably, in the one-stage training scenario, mixing code and math tokens effectively alleviates the catastrophic forgetting observed in two-stage training, while simultaneously benefiting coding performance (Table \ref{tab:code-math-general-eval}) and program-driven mathematical reasoning (Table \ref{tab:code-math}).

Code training also enhances mathematical reasoning when tools are not employed. Within the two-stage training paradigm, the initial code training phase already yields moderate improvements. It further accelerates the effectiveness of the following math training, culminating in the highest overall performance. Conversely, combining code and math tokens in one-stage training harms mathematical reasoning without tool support. A possible explanation is that DeepSeek-LLM 1.3B, due to its limited model size, cannot fully integrate both code and math data simultaneously.

\subsubsection{ArXiv Papers Appear Ineffective in Enhancing Mathematical Reasoning}

ArXiv papers are frequently incorporated as part of the mathematical pre-training datasets \citep{minerva,gpt-f,llemma,mathpile}. Nevertheless, there has been limited in-depth investigation into their actual effect on mathematical reasoning capabilities. Surprisingly, our experiments suggest that arXiv papers do not contribute effectively to improving mathematical reasoning. We conduct tests using models of varying scales, including DeepSeek-LLM 1.3B and DeepSeek-Coder-Base-v1.5 7B \citep{deepseek-coder}, trained on arXiv corpora processed through different pipelines:

\begin{itemize}[topsep=0pt]
    \item \textbf{MathPile} \citep{mathpile}: an 8.9B-token corpus curated with cleaning and filtering heuristics, where over 85\% consists of scientific arXiv papers;
    \item \textbf{ArXiv-RedPajama} \citep{redpajama}: the full set of arXiv LaTeX source files with preambles, comments, macros, and bibliographies removed, totaling 28.0B tokens.
\end{itemize}

In our trials, we train DeepSeek-LLM 1.3B for 150B tokens and DeepSeek-Coder-Base-v1.5 7B for 40B tokens separately on each arXiv dataset. The results indicate that arXiv papers do not enhance mathematical reasoning. Models trained exclusively on arXiv data show no significant gains or even exhibit declines on various mathematical benchmarks of differing difficulty levels used in this work. These benchmarks encompass quantitative reasoning datasets such as GSM8K and MATH (Table \ref{tab:arxiv-cot}), multiple-choice exams like MMLU-STEM (Table \ref{tab:arxiv-cot}), and formal mathematics tasks such as miniF2F (Table \ref{tab:arxiv_atp}).

However, this finding has certain caveats and should be interpreted cautiously. We have not yet explored:

\begin{itemize}[topsep=0pt]
    \item The influence of arXiv tokens on particular math-related tasks excluded from this study, for instance, theorem informalization, which involves transforming formal statements or proofs into informal versions;
    \item The impact of arXiv tokens when integrated with other categories of data;
    \item Whether the advantages of arXiv papers might emerge at larger model scales.
\end{itemize}

Hence, additional research is needed, which we defer to future work.

\subsection{Insights of Reinforcement Learning}
\subsubsection{Towards a Unified Paradigm} This section introduces a unified framework for analyzing various training techniques, including SFT, RFT, DPO, PPO, and GRPO, and presents experiments to investigate factors within this framework. Generally, the gradient with respect to the parameter $\theta$ for a given training method can be expressed as:

\begin{equation}
    \nabla_{\theta}\mathcal{J}_{\textcolor{red}{\mathcal{A}}}(\theta) = \mathbb{E}[\underbrace{(q,o) \sim \textcolor{red}{\mathcal{D}}}_{\text{Data Source}}]\left( \frac{1}{|o|} \sum_{t=1}^{|o|}  \underbrace{GC_{{\mathcal{A}}}(q, o, t, \textcolor{red}{\pi_{{rf}}})}_{\text{Gradient Coefficient}}  \nabla_{\theta}\log \pi_{\theta}(o_t | q, o_{<t})\right).
\label{eq:objective}
\end{equation}

There are three fundamental elements: 1) \textit{Data Source $\mathcal{D}$}, which specifies the training dataset; 2) \textit{Reward Function $\pi_{{rf}}$}, which provides the reward signals for training; 3) \textit{Algorithm $\mathcal{A}$}, which utilizes the training data and reward signals to compute the gradient coefficient $GC$, dictating the extent of punishment or reinforcement applied to the data. We examine several notable approaches within this unified framework:

\begin{itemize}
    \item \textbf{Supervised Fine-tuning (SFT)}: SFT adapts a pretrained model by training it on human-curated SFT datasets.
    \item \textbf{Rejection Sampling Fine-tuning (RFT)}: RFT further refines the SFT model by fine-tuning it using filtered outputs generated by the SFT model in response to SFT questions. The filtering is based on the accuracy of the answers.
    \item \textbf{Direct Preference Optimization (DPO)}: DPO improves the SFT model by fine-tuning it on enhanced outputs produced by the SFT model, applying a pair-wise DPO loss.
    \item \textbf{Online Rejection Sampling Fine-tuning (Online RFT)}: Unlike RFT, Online RFT starts with the SFT model as the initial policy and improves it by fine-tuning on augmented outputs sampled from the current policy model in real-time.
    \item \textbf{PPO/GRPO}: PPO/GRPO initializes the policy model from the SFT model and reinforces it through outputs sampled from the live policy model.
\end{itemize}

We present an overview of the components of these methods in Table \ref{tab:data_policy}. For a more comprehensive derivation, please see Appendix \ref{app:analysis-rl}.

\paragraph{Observation about Data Source} We categorize the data sources into two types: online sampling and offline sampling. Online sampling refers to training data generated from the exploration outcomes of the real-time training policy model, whereas offline sampling refers to training data obtained from samples generated by the initial SFT model. RFT and DPO employ the offline approach, while Online RFT and GRPO utilize the online approach.

As illustrated in Figure \ref{fig:rl-analysis}, we observe that Online RFT significantly outperforms RFT on two benchmark datasets. Specifically, Online RFT performs similarly to RFT during the early phases of training but develops a clear advantage in later stages, underscoring the benefits of online training. This is expected since, in the early stages, the actor and the SFT model are quite similar, resulting in sampled data with only slight variations. However, in later stages, the data sampled from the actor diverges more substantially, making real-time data sampling increasingly beneficial.

\paragraph{Observation about Gradient Coefficient} The algorithm transforms the reward signal into a gradient coefficient used for updating model parameters. In our experiments, we classify the reward function into two types: ‘Rule’ and ‘Model.’ ‘Rule’ refers to evaluating response quality based on answer correctness, whereas ‘Model’ involves training a reward model to assign scores to each response. The training data for the reward model is derived from rule-based judgments. Equations \ref{eq:GC-RFT} and \ref{eq:GC-GRPO} emphasize a crucial distinction between GRPO and Online RFT: GRPO uniquely modulates its gradient coefficient according to the reward values generated by the reward model. This facilitates differential reinforcement and punishment of responses depending on their reward magnitudes. In contrast, Online RFT does not incorporate this mechanism; it neither penalizes incorrect responses nor varies reinforcement among correct answers, treating all correct responses with equal intensity.

As illustrated in Figure \ref{fig:rl-analysis}, GRPO outperforms online RFT, underscoring the effectiveness of adjusting positive and negative gradient coefficients. Moreover, GRPO+PS demonstrates better results than GRPO+OS, highlighting the advantages of employing fine-grained, step-aware gradient coefficients. Additionally, we investigate iterative RL by performing two rounds of iteration in our experiments. As depicted in Figure \ref{fig:iter-rl}, iterative RL notably enhances performance, particularly during the first iteration.

\subsubsection{Why Does RL Work?} In this study, we apply reinforcement learning to a subset of instruction tuning data, achieving a substantial performance improvement over the instruction-tuned model. To further clarify the reason behind RL's effectiveness, we assess the Pass@K and Maj@K accuracy of both the Instruct and RL models on two benchmarks. As shown in Figure \ref{fig:maj-pass}, RL improves Maj@K accuracy but not Pass@K. These results suggest that RL boosts the overall model performance by making the output distribution more robust; in other words, \textbf{the improvement appears to stem from enhancing the likelihood of correct responses within the TopK outputs rather than from fundamental capability enhancements.} Similarly, \citep{wang2023making} discovered a \textbf{misalignment issue} in reasoning tasks within the SFT model, demonstrating that the reasoning abilities of SFT models can be enhanced through a series of preference alignment approaches \citep{yuan2023rrhf,song2023preference,wang2023making}.

\subsubsection{How Can RL Be Made More Effective?}
\label{sec:effec_rl}
We show that RL performs quite well in mathematical reasoning tasks and propose a unified framework to interpret various representative training methods. Within this framework, all methods are viewed as either direct or simplified RL approaches. As summarized in Equation \ref{eq:objective}, three essential elements exist: Data Source, Algorithm, and Reward Function. We discuss potential future directions concerning these components.

\paragraph{Data Source} The data source serves as the foundational material for all training methods. Specifically, in the RL context, it refers to unlabeled questions with outputs sampled from the policy model. In this work, we limit ourselves to questions from the instruction tuning phase and apply a basic nucleus sampling method to generate outputs. We believe this may explain why our RL pipeline improves only the Maj@K performance. Going forward, we plan to explore our RL framework on out-of-distribution question prompts, combined with \textbf{advanced sampling (decoding) strategies}, such as those based on tree-search methods \citep{yao2023tree}. Additionally, \textbf{efficient inference techniques} \citep{xia-etal-2023-speculative,leviathan2023fast,kwon2023efficient,xia2024unlocking}, which govern the exploration efficiency of policy models, will also play a crucial role.

\paragraph{Algorithms} Algorithms utilize the data and reward signals to calculate the gradient coefficient, which is then used to update the model parameters. According to Equation \ref{eq:objective}, all current methods to some degree completely \textbf{TRUST} the reward function's signal to either increase or decrease the conditional probability of a given token. Nonetheless, it is unrealistic to guarantee that the reward signal is always accurate, especially in highly complex tasks. For instance, even the PRM800K datasets \citep{lightman2023let}, which have been meticulously labeled by skilled annotators, still include about 20\% incorrect annotations\footnote{\url{https://github.com/openai/prm800k/issues/12\#issuecomment-1728491852}}. Therefore, we intend to investigate reinforcement learning algorithms that maintain robustness in the presence of noisy reward signals. We believe that these \textbf{WEAK-TO-STRONG} \citep{burns2023weak} alignment methods could fundamentally transform learning algorithms.

\paragraph{Reward Function} The reward function serves as the origin of the training signal. In reinforcement learning, this reward function is typically represented by a neural reward model. We identify three key research directions for reward models: 1) \textbf{Improving the generalization capability of the reward model.} It is essential for the reward model to generalize effectively to out-of-distribution queries and complex decoding outputs; without this, reinforcement learning might only stabilize the LLMs' distribution rather than enhance their core abilities; 2) \textbf{Capturing the uncertainty within the reward model.} This uncertainty could potentially bridge the gap between weak reward models and weak-to-strong learning algorithms; 3) \textbf{Developing efficient methods for constructing high-quality process reward models} that can deliver fine-grained training signals throughout the reasoning process \citep{lightman2023let,wang2023math}.

\section{Conclusion, Limitation, and Future Work} We introduce DeepSeekMath, a model that surpasses all open-source counterparts on the competition-level MATH benchmark and approaches the performance of proprietary models. DeepSeekMath~starts from DeepSeek-Coder-v1.5 7B and undergoes continuous training on 500B tokens, with a substantial portion—120B tokens—being mathematics-related data drawn from Common Crawl. Our comprehensive ablation studies reveal that web pages hold great promise for sourcing high-quality mathematical data, whereas arXiv contributions appear less advantageous than initially anticipated. We propose Group Relative Policy Optimization (GRPO), a variant of Proximal Policy Optimization (PPO), which significantly enhances mathematical reasoning while consuming less memory. Experimental outcomes demonstrate that GRPO remains effective even after DeepSeekMath-Instruct 7B attains high benchmark scores. Additionally, we offer a unified framework to interpret a sequence of approaches and highlight several promising avenues for achieving more efficient reinforcement learning.

While DeepSeekMath~achieves strong results on quantitative reasoning tasks, its performance on geometry and theorem-proving problems lags behind that of closed models. For example, during preliminary evaluations, the model struggled with questions involving triangles and ellipses, suggesting potential biases in the selection of pre-training and fine-tuning datasets. Moreover, constrained by model size, DeepSeekMath~underperforms compared to GPT-4 in few-shot learning scenarios. GPT-4 demonstrates enhanced capabilities with few-shot inputs, whereas DeepSeekMath~shows comparable performance under both zero-shot and few-shot settings. Moving forward, we plan to refine our data selection process to build a higher-quality pre-training corpus. Furthermore, we aim to investigate the promising directions outlined in Section \ref{sec:effec_rl} to develop more effective reinforcement learning techniques for large language models.

\end{document}